\chapter{Learned Image Compression}

\section{Table of Contents}

\begin{itemize}
  \item Core Idea
  \item Key Concepts
  \item Theory and Formulas
  \item Visual Schemes
  \item Examples
\end{itemize}

\section{Core Idea}

\textbf{Learned Image Compression (LIC)}, also called \textbf{Neural Image Compression (NIC)}, replaces hand-designed codec blocks with neural networks trained end-to-end on image data.

Classical codecs use fixed linear transforms:

\begin{itemize}
  \item JPEG: block DCT,
  \item JPEG 2000: DWT,
  \item handcrafted quantization and entropy models.
\end{itemize}

Learned codecs use neural networks:

\begin{itemize}
  \item \textbf{analysis transform} $g_a$: maps image to latent tensor,
  \item \textbf{quantization}: maps latents to discrete values,
  \item \textbf{entropy model}: predicts bit cost of latents,
  \item \textbf{synthesis transform} $g_s$: reconstructs the image.
\end{itemize}

\begin{calloutbox}{red!70!black}{Main paradigm shift}
Learned compression is non-linear transform coding. The transform, probability model, and reconstruction behavior are optimized from data for a rate-distortion objective instead of being manually specified.
\end{calloutbox}

\section{Key Concepts}

\subsection{Classical vs. Learned Transform Coding}

General transform-coding structure remains:

\[
x \xrightarrow{\text{analysis}} z=g_a(x)
\xrightarrow{\text{quantization}} \hat{z}=Q(z)
\xrightarrow{\text{synthesis}} \hat{x}=g_s(\hat{z})
\]

In classical codecs, $g_a$ and $g_s$ are fixed transforms such as DCT or DWT. In learned compression, $g_a$ and $g_s$ are trained CNNs.

\begin{center}\fitwidth{%
\begin{tabular}{|l|l|l|}
\hline
\textbf{Aspect} & \textbf{Classical} & \textbf{Learned} \\ \hline
Transform & fixed linear basis & learned non-linear mapping \\ \hline
Adaptivity & parameters only & weights learned from data \\ \hline
Coefficients & transform coefficients & latent variables / latent tensor \\ \hline
Optimization & discrete tool selection & gradient descent \\ \hline
Entropy model & fixed or handcrafted & learned prior/hyperprior \\ \hline
\end{tabular}}
\end{center}

\subsection{KLT as Linear Limit}

KLT is optimal among linear transforms for Gaussian data, because it decorrelates components. Neural codecs generalize this:

\begin{center}\fitwidth{%
\begin{tabular}{|l|l|}
\hline
\textbf{KLT} & \textbf{Neural codec} \\ \hline
linear & non-linear \\ \hline
optimal mainly for Gaussian sources & handles complex natural-image distributions \\ \hline
data-dependent but expensive & CNNs process large images with shared weights \\ \hline
decorrelates second-order statistics & learns high-order features: edges, textures, semantics \\ \hline
\end{tabular}}
\end{center}

\begin{calloutbox}{red!70!black}{Non-linear KLT view}
LIC can be interpreted as a non-linear KLT optimized for rate-distortion, where the basis/functions are learned instead of fixed.
\end{calloutbox}

\subsection{Rate-Distortion Objective}

Learned codecs are trained by minimizing a Lagrangian:

\[
\mathcal{L}=D(x,\hat{x})+\lambda R(\hat{z})
\]

or equivalently:

\[
\mathcal{L}=R(\hat{z})+\lambda D(x,\hat{x})
\]

The meaning of "large $\lambda$" depends on which form is used:

\begin{itemize}
  \item in $D+\lambda R$, large $\lambda$ penalizes rate more strongly;
  \item in $R+\lambda D$, large $\lambda$ penalizes distortion more strongly.
\end{itemize}

Training several models with different $\lambda$ values gives different points on the rate-distortion curve, similar in role to quality factors.

\subsection{Artificial Neuron and MLP}

Neuron:

\[
y=f(w^Tx+b)=f(w_1x_1+\cdots+w_nx_n+b)
\]

where $f$ is a non-linear activation.

Common activations:

\begin{center}\fitwidth{%
\begin{tabular}{|l|l|l|}
\hline
\textbf{Activation} & \textbf{Formula} & \textbf{Role} \\ \hline
ReLU & $f(z)=\max(0,z)$ & sparse non-linearity \\ \hline
Sigmoid & $f(z)=1/(1+e^{-z})$ & maps to $(0,1)$ \\ \hline
Tanh & $f(z)=\tanh(z)$ & maps to $(-1,1)$ \\ \hline
GDN & normalized response & compression-specific \\ \hline
\end{tabular}}
\end{center}

MLP with one hidden layer:

\[
h=f(W_1x+b_1)
\]

\[
y=g(W_2h+b_2)
\]

Non-linear activations matter because without them a deep stack collapses to one linear transform.

\subsection{Gradient Descent and Backpropagation}

Neural training updates parameters $\theta$ by:

\[
\theta_{t+1}=\theta_t-\eta\nabla_\theta\mathcal{L}(\theta_t)
\]

where $\eta$ is the learning rate.

Backpropagation applies the chain rule from output to input to compute gradients for every weight. This makes end-to-end codec training possible.

\subsection{Why MLPs Fail on Images}

Flattening an image loses spatial structure. A $256x256$ image gives:

\[
256\cdot256=65536
\]

input values. Connecting it to 1000 hidden neurons requires:

\[
65,536,000
\]

weights in one layer.

Problems:

\begin{itemize}
  \item no translation invariance,
  \item parameter explosion,
  \item overfitting,
  \item no explicit locality,
  \item computationally impractical.
\end{itemize}

\subsection{CNNs for Learned Compression}

CNNs solve image scalability using:

\begin{center}\fitwidth{%
\begin{tabular}{|l|l|}
\hline
\textbf{Feature} & \textbf{Compression meaning} \\ \hline
local connectivity & nearby pixels treated together \\ \hline
weight sharing & same detector reused everywhere \\ \hline
strided convolution & learned downsampling in encoder \\ \hline
multiple channels & many parallel learned features \\ \hline
transposed convolution & learned upsampling in decoder \\ \hline
\end{tabular}}
\end{center}

A convolution applies a small kernel, e.g. $3x3$, to every local patch. A stride $s>1$ reduces spatial resolution and compacts information.

Typical neural codec latent tensor:

\[
H\times W \rightarrow \frac{H}{16}\times\frac{W}{16}\times C
\]

after four stride-2 stages. Source notes JPEG AI-like examples with:

\begin{itemize}
  \item 160 channels for luminance,
  \item 96 channels for chrominance.
\end{itemize}

\subsection{GDN}

\begin{calloutbox}{red!70!black}{Generalized Divisive Normalization}
\[
w_i=
\frac{v_i}
{\sqrt{\beta_i+\sum_j\gamma_{ij}v_j^2}}
\]

$\beta_i$ and $\gamma_{ij}$ are learned parameters.
\end{calloutbox}

GDN is useful because it:

\begin{itemize}
  \item performs local feature normalization,
  \item mimics lateral inhibition and visual masking,
  \item reduces correlation,
  \item makes latent distributions easier to model for entropy coding.
\end{itemize}

\subsection{Autoencoder Codec}

Learned image codecs are autoencoders:

\begin{enumerate}
  \item \textbf{Analysis transform}:
\end{enumerate}

\[
y=g_a(x)
\]

\begin{enumerate}
  \item \textbf{Bottleneck}:
\end{enumerate}

\[
\hat{y}=Q(y)
\]

\begin{enumerate}
  \item \textbf{Entropy coding} of $\hat{y}$.
  \item \textbf{Synthesis transform}:
\end{enumerate}

\[
\hat{x}=g_s(\hat{y})
\]

The bottleneck is where compression happens.

\subsection{VAE Interpretation}

Modern neural codecs can be interpreted as rate-distortion VAEs.

\begin{calloutbox}{red!70!black}{VAE rate-distortion objective}
\[
\mathcal{L}
=
D_{KL}[q(\hat{y}|x)\|p(\hat{y})]
+
\lambda\mathbb{E}[\rho(x,\hat{x})]
\]

The KL term corresponds to rate, and the reconstruction term corresponds to distortion.
\end{calloutbox}

Interpretation:

\begin{itemize}
  \item $q(\hat{y}|x)$: distribution induced by encoder and quantization.
  \item $p(\hat{y})$: entropy model expected by arithmetic coder.
  \item $\rho(x,\hat{x})$: distortion metric, e.g. MSE, MS-SSIM, perceptual loss.
  \item $\lambda$: rate-distortion tradeoff parameter.
\end{itemize}

\subsection{Quantization Training Problem}

Hard rounding is non-differentiable:

\[
\hat{z}=Q(z)=\mathrm{round}(z)
\]

Its derivative is zero almost everywhere, so backpropagation cannot train the encoder through it.

Common training solution:

\[
\hat{z}\approx z+u,
\qquad
u\sim\mathcal{U}(-0.5,0.5)
\]

During inference:

\[
\hat{z}=\mathrm{round}(z)
\]

Additive uniform noise is a differentiable proxy and matches high-resolution quantization-noise theory.

Alternative approaches:

\begin{itemize}
  \item \textbf{STE}: pretend quantizer derivative is 1.
  \item \textbf{Soft quantization}: smooth approximation during training, hard rounding at test time.
\end{itemize}

\subsection{Hyperprior}

Basic latent priors assume fixed statistics, but real latent tensors have spatially varying variance. Hyperprior models transmit compact side information to predict local latent distributions.

Hyperprior mechanism:

\[
x\xrightarrow{g_a}y\xrightarrow{Q}\hat{y}\xrightarrow{g_s}\hat{x}
\]

\[
y\xrightarrow{h_a}z\xrightarrow{Q}\hat{z}\xrightarrow{h_s}\sigma
\]

Conditional model:

\[
p(\hat{y}|\hat{z})=\mathcal{N}(\mu,\sigma^2)
\]

Total rate includes:

\[
R_{\text{total}}=R(\hat{y})+R(\hat{z})
\]

The hyperprior overhead is useful when the rate saved by a better model for $\hat{y}$ exceeds the rate spent on $\hat{z}$.

\subsection{JPEG AI}

Source notes describe \textbf{JPEG AI} as an international learned-image-coding standard built on deep learning. Its goals:

\begin{itemize}
  \item outperform HEVC/VVC intra coding in image compression efficiency,
  \item support both human and machine tasks,
  \item use hierarchical VAE-like neural codecs,
  \item provide complexity profiles for different hardware.
\end{itemize}

Key ideas:

\begin{itemize}
  \item deep CNN analysis/synthesis transforms,
  \item hyperprior side information,
  \item quantization proxy during training,
  \item latent tensors with many channels,
  \item one bitstream usable by human reconstruction or machine task decoders.
\end{itemize}

Complexity profiles in source:

\begin{center}\fitwidth{%
\begin{tabular}{|l|l|l|}
\hline
\textbf{Profile} & \textbf{Complexity} & \textbf{Target} \\ \hline
Dec0 & 8 kMAC/px & low-end CPUs \\ \hline
Dec1 & 23 kMAC/px & real-time mid-range smartphones \\ \hline
Dec2 & 214 kMAC/px & high-end GPU / maximum quality \\ \hline
\end{tabular}}
\end{center}

\subsection{Current Limitations}

\begin{center}\fitwidth{%
\begin{tabular}{|l|l|}
\hline
\textbf{Limitation} & \textbf{Why it matters} \\ \hline
high complexity & many kMAC/px, high energy cost \\ \hline
hardware dependence & GPU/NPU often needed \\ \hline
determinism & floating-point differences can cause drift \\ \hline
out-of-distribution risk & unusual images may be poorly reconstructed \\ \hline
hallucination & neural decoder may synthesize plausible but false details \\ \hline
\end{tabular}}
\end{center}

\subsection{Future Directions}

\begin{itemize}
  \item \textbf{Generative compression:} GAN/diffusion decoders synthesize realistic textures at very low rates.
  \item \textbf{Coding for machines:} transmit latents directly to detection/segmentation systems.
  \item \textbf{Hardware-algorithm co-design:} neural codecs designed with NPUs and mobile accelerators.
\end{itemize}

\section{Theory and Formulas}

\subsection{Classical-to-Neural Architecture}

Classical transform coding:

\[
x\rightarrow T(x)\rightarrow Q(T(x))\rightarrow \text{entropy coding}
\]

Learned transform coding:

\[
x\rightarrow g_a(x)=y\rightarrow Q(y)=\hat{y}
\rightarrow \text{entropy coding}
\rightarrow g_s(\hat{y})=\hat{x}
\]

The learned transform is non-linear, data-driven, and optimized jointly with the entropy model.

\subsection{Latent Tensor Size}

If an input image has shape $H\times W$, and the encoder uses four stride-2 stages:

\[
H\times W
\rightarrow
\frac{H}{2}\times\frac{W}{2}
\rightarrow
\frac{H}{4}\times\frac{W}{4}
\rightarrow
\frac{H}{8}\times\frac{W}{8}
\rightarrow
\frac{H}{16}\times\frac{W}{16}
\]

With $C$ channels:

\[
y\in\mathbb{R}^{H/16\times W/16\times C}
\]

Spatial size shrinks, but channel count increases to preserve multiple learned features.

\subsection{Entropy Model and Arithmetic Coding}

If entropy model assigns probability $p_{\hat{y}}(\hat{y}_i)$ to each quantized latent symbol, rate estimate is:

\[
R(\hat{y})\approx
-\sum_i \log_2 p_{\hat{y}}(\hat{y}_i)
\]

Better probability models reduce arithmetic-coding length.

With hyperprior:

\[
R(\hat{y}|\hat{z})
\approx
-\sum_i\log_2 p(\hat{y}_i|\hat{z})
\]

Total:

\[
R_{\text{total}}
=
R(\hat{z})+R(\hat{y}|\hat{z})
\]

\subsection{Additive Noise Approximation}

Uniform scalar quantization with unit step has error approximately:

\[
e\sim\mathcal{U}(-0.5,0.5)
\]

So during training:

\[
\tilde{y}=y+u,\qquad u\sim\mathcal{U}(-0.5,0.5)
\]

approximates:

\[
\hat{y}=\mathrm{round}(y)
\]

while preserving gradient flow.

\subsection{GDN Interpretation}

GDN divides each feature by local energy:

\[
w_i=\frac{v_i}{\sqrt{\beta_i+\sum_j\gamma_{ij}v_j^2}}
\]

If neighboring channels are strong, $w_i$ is suppressed. This resembles perceptual masking and helps produce latents with simpler distributions.

\section{Visual Schemes}

\subsection{Neural Network Basics}

\begin{center}\includegraphics[width=\linewidth]{ch06/artificial-neuron.png}\end{center}
Artificial neuron computes weighted sum, adds bias, and applies non-linear activation.

\begin{center}\includegraphics[width=0.833\linewidth]{ch06/mlp-architecture.png}\end{center}
MLP connects all neurons between layers; useful conceptually but inefficient for images.

\begin{center}\includegraphics[width=0.833\linewidth]{ch06/image-flattening-problem.png}\end{center}
Flattening image pixels destroys spatial topology and causes parameter explosion.

\subsection{CNN Features}

\begin{center}\includegraphics[width=\linewidth]{ch06/convolution-receptive-field.png}\end{center}
Convolution applies the same local kernel at all positions, reusing learned pattern detectors.

\begin{center}\includegraphics[width=\linewidth]{ch06/strided-convolution.png}\end{center}
Strided convolution downsamples spatial dimensions and compacts information.

\begin{center}\includegraphics[width=\linewidth]{ch06/parallel-feature-channels.png}\end{center}
Multiple kernels create multiple feature maps, forming a 3D latent representation.

\subsection{Autoencoder Codec}

\begin{center}\footnotesize\itshape Diagram (Mermaid source):\end{center}
\begin{verbatim}
flowchart LR
    A[Input image x] --> B[Analysis transform ga]
    B --> C[Latent tensor y]
    C --> D[Quantization Q]
    D --> E[Quantized latents yhat]
    E --> F[Entropy coding]
    F --> G[Bitstream]
    G --> H[Entropy decoding]
    H --> E2[Quantized latents yhat]
    E2 --> I[Synthesis transform gs]
    I --> J[Reconstructed image xhat]
\end{verbatim}

\subsection{Hyperprior Codec}

\begin{center}\footnotesize\itshape Diagram (Mermaid source):\end{center}
\begin{verbatim}
flowchart LR
    A[Image x] --> B[Analysis ga]
    B --> Y[Latents y]
    Y --> QY[Quantize yhat]
    Y --> HA[Hyper-analysis ha]
    HA --> Z[Hyperlatents z]
    Z --> QZ[Quantize zhat]
    QZ --> HS[Hyper-synthesis hs]
    HS --> P[Scale / probability model]
    QY --> EC[Arithmetic coding]
    P --> EC
    QZ --> ECZ[Encode side info]
    EC --> BS[Bitstream]
    ECZ --> BS
    BS --> DEC[Decode yhat, zhat]
    DEC --> GS[Synthesis gs]
    GS --> XH[Image xhat]
\end{verbatim}

\subsection{JPEG AI Dual Use}

\begin{center}\footnotesize\itshape Diagram (Mermaid source):\end{center}
\begin{verbatim}
flowchart LR
    A[JPEG AI bitstream] --> B[Decode latent tensor]
    B --> C[Human vision decoder]
    B --> D[Machine task decoder]
    C --> E[RGB reconstruction]
    D --> F[Detection / segmentation / features]
\end{verbatim}

\begin{calloutbox}{blue!60!black}{Missing visual}
Source has real figures for neural-network/CNN basics, but not for autoencoder, hyperprior, or dual-use JPEG AI architecture. Mermaid schemes above summarize those architectures.
\end{calloutbox}

\section{Examples}

\begin{calloutbox}{green!55!black}{MLP parameter explosion}
A $256x256$ image has:

\[
256\cdot256=65536
\]

input values. A fully connected layer with 1000 neurons needs:

\[
65536\cdot1000=65,536,000
\]

weights, before biases.

\textbf{Takeaway:} CNNs are needed because they reuse local filters instead of connecting every pixel to every neuron.
\end{calloutbox}

\begin{calloutbox}{green!55!black}{Latent tensor dimensions}
With four stride-2 convolution layers:

\[
H\times W\rightarrow H/16\times W/16
\]

If $C=160$ channels, the encoder produces:

\[
y\in\mathbb{R}^{H/16\times W/16\times160}
\]

\textbf{Takeaway:} spatial resolution shrinks, but many channels store learned features.
\end{calloutbox}

\begin{calloutbox}{green!55!black}{Training vs. inference quantization}
During training:

\[
\tilde{y}=y+u,\qquad u\sim\mathcal{U}(-0.5,0.5)
\]

During inference:

\[
\hat{y}=\mathrm{round}(y)
\]

\textbf{Takeaway:} additive noise makes training differentiable while approximating quantization noise.
\end{calloutbox}

\begin{calloutbox}{green!55!black}{Hyperprior rate tradeoff}
Sending hyperlatents $\hat{z}$ costs bits, but improves probability estimates for $\hat{y}$.

Useful condition:

\[
R(\hat{z})+R(\hat{y}|\hat{z}) < R(\hat{y})
\]

\textbf{Takeaway:} side information is worth sending only when it saves more latent bits than it costs.
\end{calloutbox}

\begin{calloutbox}{green!55!black}{Complexity profiles}
\begin{center}\fitwidth{%
\begin{tabular}{|l|l|l|}
\hline
\textbf{Profile} & \textbf{Complexity} & \textbf{Interpretation} \\ \hline
Dec0 & 8 kMAC/px & lightweight decoder \\ \hline
Dec1 & 23 kMAC/px & phone-class real-time target \\ \hline
Dec2 & 214 kMAC/px & high-quality high-compute decoder \\ \hline
\end{tabular}}
\end{center}

\textbf{Takeaway:} same bitstream can be decoded with different compute/quality profiles.
\end{calloutbox}

\begin{calloutbox}{green!55!black}{Classical vs. learned codec}
\begin{center}\fitwidth{%
\begin{tabular}{|l|l|l|}
\hline
\textbf{Aspect} & \textbf{JPEG / JPEG 2000} & \textbf{Learned / JPEG AI} \\ \hline
Transform & DCT / DWT & CNN + GDN \\ \hline
Basis & hand-crafted & learned from data \\ \hline
Optimization & module-wise & end-to-end \\ \hline
Entropy model & Huffman/arithmetic tables & learned prior/hyperprior \\ \hline
Artifacts & blocking or blur & possible hallucination \\ \hline
Complexity & low & high \\ \hline
\end{tabular}}
\end{center}

\textbf{Takeaway:} learned codecs win rate-distortion performance, but cost much more computation and must manage determinism.
\end{calloutbox}
